\section{Conclusion} % (~5%)
% A conclusion may review the main points of the work, do not replicate the abstract as the conclusion. A conclusion might elaborate on the importance of the work or suggest applications, extensions, and improvements.

\par In this project, we successfully constructed a~soft robot capable of non-verbal communication through expressive movements. The~expressed emotions are basic but distinctive, even with the~encountered problems typical for the~prototyping phase.

\par The~conducted HRI experiment enabled to evaluate the~robot's design. Although the~results did not support our hypothesis regarding the~impact of ear movements on emotional expressiveness, they provided valuable data for future improvements.

\par In the~future, the~robot's electronics and power source should be enhanced. The~more powerful power source would allow for smoother movements without jittering. Switching the~motor shield for a~more advanced one, would open possibilities for adding more actuators, pumps or valves. Therefore, more complex emotions could be implemented, potentially improving the~robot's expressiveness and user engagement.

\par The~HRI experiment could be then repeated with an~improved robot design and a~larger, more diverse participant pool. Then, it could be compared with the~current results, to find out if the~new changes had a~positive impact in the~eyes of the~users.
